\centerline{\bf \Huge Смертельная битва}

\begin{flushright}
        \scriptsize{Мы, люди, хрупкие и смертные. \\Однако, даже когда нас ранят или мучат,\\ в нас всё равно не угасает жажда жизни.\\Гатс}
\end{flushright}


\textbf{Правила смертельной битвы:}

1) Назад пути нет

2) За каждого поверженного врага тройная награда

3) После того как враг повержен хотя бы одним участником, остальные не имеют права с ним биться

4) Перечитать первое правило

\problem{1}
Докажите, что если $r$ это простой делитель числа $\frac{x^p-1}{x-1}$, то либо $r \equiv 1(mod\  p)$ либо $r=p$.

\problem{2}
На клетчатой плоскости назовём бриллиантовым расстоянием между двумя точками $A$ и $B$ находящихся в узлах, количество клеток, которые пересекает отрезок $A B$. Так же назовём бриллиантовым кругом радиуса $n$ геометрическое место точек удалённых от данной на бриллиантовое расстояние $n$. Докажите, что бриллиантовый круг радиуса $n$ совпадает со множеством узлов сетки заданных уравнением $|x|+|y|=n+1$, где $x, y$ целые тогда и только тогда, когда $n+1$ простое.

\problem{3}
Назовём \textit{паспортом} натурального числа набор простых чисел, входящих в его разложение. Докажите, что в любой натуральной арифметической прогресии есть сколько угодно много чисел с одинаковым паспортом.

\problem{4}
Докажите, что для любых натуральных $x, y$ верно, что $4xy-1$ является делителем $(4x^2-1)^2$ тогда и только тогда, когда $x = y.$

\problem{5}
Пусть $P(n)$ это сумма всех простых чисел меньших $n$. Докажите, что функция $D(n) = \text{НОД}(P(n), n)$ бексонечно много раз принимает значение 1.

\problem{6}
Агаларов Агалар загадал число $n$ с суммой цифр 2025. Тралалело Тралала может за один плюсик узнать сумму цифр числа $n-m$, где $m$ он выбирает сам. Какое минимальное количество плюсиков получит Агалар прежде чем Тралалело Тралала разгадает его число?